% This LaTeX document needs to be compiled with XeLaTeX.
\documentclass[
  10
]{article}
\usepackage[margin=2cm,includehead=true,includefoot=true,centering,]{geometry}
\usepackage{xcolor}
\usepackage{ucharclasses}
\usepackage{hyperref}
\usepackage{amsmath,amssymb}
\usepackage{amsfonts}
\usepackage[version=4]{mhchem}
\usepackage{stmaryrd}
\usepackage{bbold}
\usepackage{polyglossia}
\usepackage{fontspec}
\usepackage{ctex}
\usepackage[export]{adjustbox}
\usepackage{tabularx}
\usepackage{booktabs}

\usepackage{setspace}
\setstretch{1.2}

\makeatletter
\@ifundefined{KOMAClassName}{% if non-KOMA class
  \IfFileExists{parskip.sty}{%
    \usepackage{parskip}
  }{% else
    \setlength{\parindent}{0pt}
    \setlength{\parskip}{6pt plus 2pt minus 1pt}}
}{% if KOMA class
  \KOMAoptions{parskip=half}}
\makeatother
\usepackage{longtable,booktabs,array}
\newcounter{none} % for unnumbered tables
\usepackage{multirow}
\usepackage{calc} % for calculating minipage widths
% Correct order of tables after \paragraph or \subparagraph
\usepackage{etoolbox}
\makeatletter
\patchcmd\longtable{\par}{\if@noskipsec\mbox{}\fi\par}{}{}
\makeatother
% Allow footnotes in longtable head/foot
\IfFileExists{footnotehyper.sty}{\usepackage{footnotehyper}}{\usepackage{footnote}}
\makesavenoteenv{longtable}
\usepackage{graphicx}
\makeatletter
\newsavebox\pandoc@box
\newcommand*\pandocbounded[1]{% scales image to fit in text height/width
  \sbox\pandoc@box{#1}%
  \Gscale@div\@tempa{\textheight}{\dimexpr\ht\pandoc@box+\dp\pandoc@box\relax}%
  \Gscale@div\@tempb{\linewidth}{\wd\pandoc@box}%
  \ifdim\@tempb\p@<\@tempa\p@\let\@tempa\@tempb\fi% select the smaller of both
  \ifdim\@tempa\p@<\p@\scalebox{\@tempa}{\usebox\pandoc@box}%
  \else\usebox{\pandoc@box}%
  \fi%
}
% Set default figure placement to htbp
\def\fps@figure{htbp}
\makeatother
\setlength{\emergencystretch}{3em} % prevent overfull lines
\providecommand{\tightlist}{%
  \setlength{\itemsep}{0pt}\setlength{\parskip}{0pt}}


\hypersetup{colorlinks=true, linkcolor=blue, filecolor=magenta, urlcolor=cyan,}
\urlstyle{same}

\usepackage{colortbl}
\definecolor{table-row-color}{HTML}{999999}
\definecolor{table-rule-color}{HTML}{999999}
%\arrayrulecolor{black!40}
\arrayrulecolor{table-rule-color}     % color of \toprule, \midrule, \bottomrule
\setlength{\aboverulesep}{0pt}
\setlength{\belowrulesep}{0pt}


\setotherlanguages{english}
\IfFontExistsTF{Source Han Serif CN}
{\newfontfamily\chinesefont{Source Han Serif CN}}
{\IfFontExistsTF{Noto Serif CJK SC}
  {\newfontfamily\chinesefont{Noto Serif CJK SC}}
  {\IfFontExistsTF{SimSun}
    {\newfontfamily\chinesefont{SimSun}}
    {\IfFontExistsTF{FangSong}
      {\newfontfamily\chinesefont{FangSong}}
      {\newfontfamily\chinesefont{Arial Unicode MS}}
}}}
\IfFontExistsTF{Times New Roman}
{\newfontfamily\englishfont{Times New Roman}}
{\IfFontExistsTF{Liberation Serif}
  {\newfontfamily\englishfont{Liberation Serif}}
  {\IfFontExistsTF{DejaVu Serif}
    {\newfontfamily\englishfont{DejaVu Serif}}
    {\newfontfamily\englishfont{Arial}}
}}


\author{}
\date{}

\begin{document}


\section{Seq-NMS for Video Object
Detection}\label{seq-nms-for-video-object-detection}

Wei Han1∗, Pooya Khorrami1∗, Tom Le Paine1∗, Prajit
Ramachandran1,Mohammad Babaeizadeh1, Honghui Shi1, Jianan Li2Shuicheng
Yan2, Thomas S. Huang1

1University of Illinois at Urbana-Champaign\{weihan3, pkhorra2, paine1,
prmchnd2mb2, hshi10, t-huang1\}@illinois.edu2National University of
Singapore\{elev373, eleyans\}@nus.edu.sg

\section{Abstract}\label{abstract}

Video object detection is challenging because objects that are easily
detected inone frame may be difficult to detect in another frame within
the same clip. Re-cently, there have been major advances for doing
object detection in a single im-age. These methods typically contain
three phases: (i) object proposal generation(ii) object classification
and (iii) post-processing. We propose a modification ofthe
post-processing phase that uses high-scoring object detections from
nearbyframes to boost scores of weaker detections within the same clip.
We show thatour method obtains superior results to state-of-the-art
single image object detec-tion techniques. Our method placed
\(3 ^ { r d }\) in the video object detection (VID) taskof the ImageNet
Large Scale Visual Recognition Challenge 2015 (ILSVRC2015).

\section{1 Introduction}\label{introduction}

Single image object detection has experienced large performance gains in
the last few years. Videoobject detection, on the other hand, still
remains an open problem. This is mainly because objectsthat are easily
detected in one frame may be difficult to detect in another frame within
the same videoclip. There are many reasons that may cause this
difficulty. Some examples include: (i) drastic scalechanges (ii)
occlusion and (iii) motion blur. In this work we propose a simple
extension of singleimage object detection to help overcome these
difficulties.

Current state-of-the-art single image object detection systems can be
broken up into three distinctphases: (i) region proposal generation (ii)
object classification and (iii) post-processing. During theregion
proposal generation phase, a set of candidate regions are generated
based on how likely theyare to contain an object. Previous region
proposal methods were based on low-level image features{[}11, 16{]}
while the current state-of-the-art, Faster R-CNN, {[}9{]} learns to
generate proposals using aneural network. The candidate regions are then
assigned a class score in the object classificationphase, and redundant
detections are subsequently filtered in the post-processing phase.

While effective, single image methods are na¨ıve because they completely
ignore the temporal dimen-sion. In this work, we incorporate temporal
information during the post-processing phase in order torefine the
detections within each individual frame. Given a video sequence of
region proposals andtheir corresponding class scores, our method
associates bounding boxes in adjacent frames using asimple overlap
criterion. It then selects boxes to maximize a sequence score. Those
boxes are thenused suppress overlapping boxes in their respective frames
and are subsequently rescored in orderto boost weaker detections.

\pandocbounded{\includegraphics[keepaspectratio,alt={image}]{images/f9f2109201d160e0ca3f7974293bb136803b87c3ab48f5344f60768072eb9a1c.jpg}}

Figure 1: Illustration of Seq-NMS. Seq-NMS takes as input all object
proposals boxes B and scoresS for an entire video clip V (in contrast to
NMS which takes proposals from a single image). It isapplied
iteratively. At each iteration it performs three steps: 1) Sequence
Selection, which selectsthe sequence of boxes with the highest sequence
score, \(\mathbf { B } ^ { s e q }\) . 2) Sequence Re-scoring, which
takesall scores in the sequence
\(\mathbf { S } ^ { s e q ^ { \prime } }\) and applies a function to
them to get a new score for each frame inthe sequence
\(\mathbf { S } ^ { s e q }\) . 3) Suppression, which for each box in
\(\mathbf { B } ^ { s e q }\) , suppresses any boxes in the sameframe
that have sufficient overlap.

The main contributions of our work are as follows:

\begin{enumerate}
\def\labelenumi{\arabic{enumi}.}
\item
  We present Seq-NMS, a method to improve object detection pipelines for
  video data.Specifically, we modify the post-processing phase to use
  high-scoring object detectionsfrom nearby frames in order to boost
  scores of weaker detections within the same clip.
\item
  We evaluate Seq-NMS on the ImageNet VID dataset and show that it
  outperforms state-of-the-art single image-based methods. We show that
  our method is helpful in cases wheresingle frames contain objects that
  are at extreme scales, occluded, or blurred. We presentspecific
  instances where our Seq-NMS improves performance.
\item
  Our method placed \(3 ^ { r d }\) in the video object detection (VID)
  task of the ImageNet LargeScale Visual Recognition Challenge 2015
  (ILSVRC2015).
\end{enumerate}

\section{2 Our Approach}\label{our-approach}

\section{2.1 Seq-NMS}\label{seq-nms}

Most object detection methods (Faster R-CNN included) are designed for
performing object de-tection on a single independent frame. However,
since we are concerned with object detection invideos, it would be a
waste of salient information to ignore the temporal component entirely.
Oneproblem we noticed with Faster R-CNN on the validation set was that
non-maximum suppression(NMS) frequently chose the wrong bounding box
after object classification. It would choose boxesthat were overly
large, resulting in a smaller intersection-over-union (IoU) with the
ground truth boxbecause the union of areas was large. The large boxes
often had very high object scores, possiblybecause more information is
available to be extracted during RoI pooling. In order to combat
thisproblem, we attempted to use temporal information to re-rank boxes.
We assume that neighboringframes should have similar objects, and their
bounding boxes should be similar in position and size,i.e.~temporal
consistency.

To make use of this temporal consistency, we propose a heuristic method
for re-ranking boundingboxes in video sequences called Seq-NMS. Seq-NMS
has three steps: Step 1) Sequence Selection,Step 2) Sequence Re-scoring,
Step 3) Suppression. We repeat these three steps until a no sequencesare
left. Figure 1 illustrates this process.

\pandocbounded{\includegraphics[keepaspectratio,alt={image}]{images/9dcb1a280d5f8f66ef1a5685583738fb7dc1b219f408395907a4591b6df141c3.jpg}}

\pandocbounded{\includegraphics[keepaspectratio,alt={image}]{images/876c8e04ca7a888941efd23754a176b281d0ce14348820ee42dbf7a5ecfe96a8.jpg}}

\pandocbounded{\includegraphics[keepaspectratio,alt={image}]{images/fae1bc492a681546e19e84c93b84485074b0fa09619c60dbd91d64d18dd17e32.jpg}}

Figure 2: Illustration of Sequence Selection. We construct a graph where
boxes in adjacent framesare linked iff their
\(\mathrm { I o U } > 0 . 5\) . A sequence is a set of boxes that are
linked in the video. We then selectthe sequence with the highest
sequence score shown in Equation 1. This produces
\(\mathbf { B } ^ { s e q }\) and
\(\mathbf { S } ^ { s e q ^ { \prime } }\)which is a set of at most one
box per frame, and the associated scores. After Sequence Selection,
foreach box in \(\mathbf { B } ^ { s e q }\) , we suppress any boxes in
the same frame that have \(\mathrm { I o U } > 0 . 3\) .

Seq-NMS is performed on a single video clip V which is comprised of a
set of \(T\) frames,\(\{ v _ { 0 } , \ldots , v _ { T } \}\) . For each
frame \(t\) , we have a set of region proposal boxes \(b _ { t }\) and
scores \(s _ { t }\) bothof size \(n _ { t }\) , which varies for each
frame. The set of proposals for an entire clip is denoted
by\(\mathbf { B } = \{ b _ { 0 } , \dots , b _ { T } \}\) . Likewise,
the set of scores for the entire clip is denoted by
\(\mathbf { S } = \{ s _ { 0 } , \ldots , s _ { T } \}\) .

Given a set of region bounding boxes B, and their detection scores S as
input, sequence selectionchooses a subset of boxes
\(\mathbf { B } ^ { s e q }\) and their associated scores
\(\mathbf { S } ^ { s e q ^ { \prime } }\) . The re-scoring function
takes \(\mathbf { S } ^ { s e q ^ { \prime } }\)and produces a new set
of scores \(\mathbf { S } ^ { s e q }\) .

Sequence Selection. For each pair of neighboring frames in V, a bounding
box in the first framecan be linked with a bounding box in the second
frame iff their IoU is above some threshold. Wefind potential linkages
in each pair of neighboring frames across the entire clip. Then, we
attempt tofind the maximum score sequence across the entire clip. That
is, we attempt to find the sequence ofboxes that maximize the sum of
object scores subject to the constraint that all adjacent boxes mustbe
linked.

\[
i ^ {\prime} \quad = \quad \underset {i _ {t _ {s}}, \dots , i _ {t _ {e}}} {\operatorname {a r g m a x}} \sum_ {t = t _ {s}} ^ {t _ {e}} s _ {t} [ i _ {t} ]
\]

\begin{enumerate}
\def\labelenumi{(\arabic{enumi})}
\tightlist
\item
\end{enumerate}

This can be found efficiently using a simple dynamic programming
algorithm that maintains themaximum score sequence so far at each box.
The optimization returns a set of indices \(i ^ { \prime }\) thatare
used to extract a sequence of boxes
\(\mathbf { B } ^ { s e q } = \{ b _ { t _ { s } } [ i _ { t _ { s } } ] , . . . , b _ { t _ { e } } [ i _ { t _ { e } } ] \}\)
and their scores
\(\begin{array} { r l } { \mathbf { S } ^ { s e q ^ { \prime } } = } \end{array}\)\(\{ s _ { t _ { s } } [ i _ { t _ { s } } ] , \ldots , s _ { t _ { e } } [ i _ { t _ { e } } ] \}\)
. Figure 2 gives a visual example of the sequence selection phase.

Sequence Re-scoring. After the sequence is selected, the scores within
it are improved. We applya function \(F\) to the sequence scores to
produce
\(\mathbf { S } ^ { s e q } = F ( \mathbf { S } ^ { s e q ^ { \prime } } )\)
. We try two different re-scoringfunctions: the average and the max.

Suppression. The boxes in the sequence are then removed from the set of
boxes we link over.Furthermore, we apply suppression within frames such
that if a bounding box in frame \(t , t \in [ t _ { s } , t _ { e } ]\)
,has an IoU with \(b _ { t }\) over some threshold, it is also removed
from the set of candidate boxes.

Table 1: Number of Samples in Imagenet VID Dataset

{\def\LTcaptype{none} % do not increment counter
\begin{longtable}[]{@{}lllll@{}}
\toprule\noalign{}
\endhead
\bottomrule\noalign{}
\endlastfoot
\multicolumn{2}{@{}l}{%
} & Train & Validation & Test \\
\multirow{2}{*}{Initial} & Snippets & 1,952 & 281 & 458 \\
& Images & 405,014 & 64,698 & 127,618 \\
\multirow{2}{*}{Full} & Snippets & 3,862 & 555 & 937 \\
& Images & 1,122,397 & 176,126 & 315,176 \\
\end{longtable}
}

\section{3 The Dataset}\label{the-dataset}

For the 2015 iteration, the ImageNet competition contained a new taster
competition for object de-tection from video called the ImageNet VID
competition. Similar to the ImageNet object detectiontask (DET), the
task is to classify and locate objects in every image. However, instead
of containinga collection of independent images, the VID dataset groups
several frames from the same video intovideo clips or ''snippets''. All
visible objects in every frame are annotated with a class label
andbounding box. The VID dataset contains 30 object categories which are
a subset of the 200 cate-gories provided in the DET dataset. The dataset
contains three sets of non-overlapping videos andlabels: train,
validation, and test. The training, validation and test sets in the
initial release of theVID dataset contain 1,952, 281 and 458 snippets
respectively. Meanwhile, the final release roughlydoubled the number
snippets in each set to 3,862, 555, and 937. The number of snippets and
numberof images in each set of the ImageNet VID dataset can be found in
Table 1.

\section{4 Results}\label{results}

\section{4.1 Training Details for RPN and
Classifier}\label{training-details-for-rpn-and-classifier}

In Faster R-CNN, the RPN and the classification network share
convolutional layers and are trainedtogether in an alternating fashion.
First, we trained a Zeiler Fergus (ZF) style {[}14{]} RPN
usingstochastic gradient descent and the image sampling strategy
described in {[}5{]}. We accomplished thisby first training the RPN on
the initial VID training dataset for 400,000 iterations. We then
traineda ZF style Fast R-CNN on the initial VID training set for 200,000
iterations. Finally, we refined theRPN by fixing the convolutional
layers to be those of the trained detector and trained for 400,000steps.
We found that our trained RPN was able to obtain proposals that
overlapped with the groundtruth boxes in the initial VID validation set
with recall over \(90 \%\) .

For our classifier, we considered both a Zeiler Fergus style network (ZF
net) and VGG16 network(VGG net) {[}10{]}. The ZF network was trained on
the initial VID training set and the VGG16 netwas pre-trained on the
training and validation sets of the 2015 ImageNet DET challenge. The
DETdataset contained 200 object categories and the train and validation
sets contained 456,567 and55,502 images, respectively. We then replaced
the 200 unit softmax layer with a 30 unit one andtrained it on the
initial VID training set (405K images) while keeping all of the other
layers fixed.It should be noted that we never used the full training set
(1.1M images) in any of our experiments.Our models were trained using a
heavily modified version of the open source Faster R-CNN Caffecode
released by the authors1.

\section{4.2 Quantitative Results}\label{quantitative-results}

We validated our method by conducting experiments on the initial and
full validation set as well asthe full test set of the ImageNet VID
dataset. During the post-processing phase, we considered fourdifferent
techniques: (i) single image NMS (ii) Seq-NMS (avg) (iii) Seq-NMS (max)
(iv) Seq-NMS(best). Seq-NMS (avg) and Seq-NMS (max) rescored the
sequences selected by Seq-NMS using theaverage or max detection scores
respectively, while Seq-NMS (best) chose the best performing ofthe three
aforementioned techniques on each class and averaged the results.

Table 2 shows our results on the initial and full validation set. We
found that using VGG net gavea substantial improvement over using the
architecture described by Zeiler and Fergus. Sequence

Table 2: Method comparison on initial and full ImageNet VID validation
set

{\def\LTcaptype{none} % do not increment counter
\begin{longtable}[]{@{}|l|l|l|@{}}
\toprule\noalign{}
\endhead
\bottomrule\noalign{}
\endlastfoot
\hline
Method & mAP(\%) - (Initial Val) & mAP(\%) - (Full Val) \\
\hline
ZF net + NMS & 32.2 & - \\
\hline
ZF net + Seq-NMS (max) & 36.3 & - \\
\hline
ZF net + Seq-NMS (avg) & 38.3 & - \\
\hline
ZF net + Seq-NMS (best) & 40.2 & - \\
\hline
VGG net + NMS & 44.4 & 44.9 \\
\hline
VGG net + Seq-NMS (max) & 50.1 & 50.5 \\
\hline
VGG net + Seq-NMS (avg) & 51.5 & 51.4 \\
\hline
VGG net + Seq-NMS (best) & 53.6 & 52.2 \\
\hline
CUVideo - T-CNN {[}7{]} & - & 73.8 \\
\hline
\end{longtable}
}

re-scoring with Seq-NMS gave further improvements. On the initial
validation set, Seq-NMS (avg)achieved a mAP score of \(5 1 . 5 \%\) .
This result can be further improved to \(5 3 . 6 \%\) when combining
allthree NMS techniques. Meanwhile on the full validation set, Seq-NMS
(avg) got a mAP score of\(5 1 . 4 \%\) . When combining all three NMS
methods (Seq-NMS (best)) on the full val set, we achieve amAP score of
\(5 2 . 2 \%\) . In Figure 3, we give a full breakdown of Seq-NMS' (avg)
performance acrossall 30 classes and compare it with the single image
NMS technique. Figure 4 shows which classesexperienced the largest gains
in performance when switching from single image NMS to Seq-NMS(avg). The
5 classes that experienced the highest gains in performance were: (i)
motorcycle (ii)turtle (iii) red panda (iv) lizard and (v) sheep.

On the test set, we ranked \(3 ^ { \mathrm { r d } }\) in terms of
overall mean average precision (mAP). The results ofVGG net models are
shown in Table 3. Once again, we see that Seq-NMS and Rescoring
showedsignificant improvements over traditional frame-wise NMS
post-processing. Our best submissionachieved a mAP of
\(4 8 . 7 \% ^ { 2 }\) .

We also report the mAP score of the challenge's top performing method
{[}7{]} on both the validationand test sets in Tables 2 and 3. In
{[}7{]}, the authors present a suite of techniques including (i) a
strongstill-image detector (ii) bounding box suppression and propagation
(iii) trajectory / tubelet re-scoringand (iv) model combination. The
still-image detector's performance achieves a mAP of \(6 7 . 7 \%\)
.When directly comparing the amount of improvement obtained just from
temporal information (ii)and (iii), our method is superior (
\(7 . 3 \%\) vs.~\(6 . 7 \%\) ).

\section{4.3 Qualitative Results}\label{qualitative-results}

In Figure 5, we present clips from the ImageNet VID dataset where
Seq-NMS improved perfor-mance. The boxes represent a sequence selected
by Seq-NMS. Clips were subsampled to provideexamples of high and low
scoring boxes. In each of these clips, the object of interest is
subjectedto one or more perturbations commonly seen in video data such
as occlusion (clips a, b, and e),drastic scaling (clip c), and blur
(clip d). These perturbations naturally cause the classifier to
scoreproposals with much lower confidence. However, since the Seq-NMS
has associated these lowerconfidence detections with previous higher
confidence detections of the same object, rescoring thelower confidence
detections with the average improves precision.

We also present instances where Seq-NMS does not appear to improve
performance in Figure 6.One case where Seq-NMS may not help is when
there are several objects with similar appearanceclose together in the
video (clip a). This will cause the detector to drift from one object to
anotherwhich leads to missed detections and incorrect score assignment.
Another case is when Seq-NMSaccumulates spurious detections which leads
to many more false positives (clips b and c). Thisoccurs because
Seq-NMS' objective function, the sum of a sequence's confidence scores,
does notpenalize against adding more detections.

Table 3: Method comparison on full ImageNet VID test set

{\def\LTcaptype{none} % do not increment counter
\begin{longtable}[]{@{}|l|l|@{}}
\toprule\noalign{}
\endhead
\bottomrule\noalign{}
\endlastfoot
\hline
Method & mAP (\%) \\
\hline
VGG net + NMS & 43.4 \\
\hline
VGG net + Seq-NMS (max) & 47.5 \\
\hline
VGG net + Seq-NMS (avg) & 48.7 \\
\hline
VGG net + Seq-NMS (best) & 48.2 \\
\hline
CUVideo - T-CNN {[}7{]} & 67.8 \\
\hline
\end{longtable}
}

\pandocbounded{\includegraphics[keepaspectratio,alt={image}]{images/00604adaac2e87fd50b474b784a124e78f2dca9f1b5fc521d9d5acc2e1b80d39.jpg}}

Figure 3: Performance (mAP) of our Seq-NMS and NMS. Performance is
measured on the full Im-ageNet validation set. We use average rescoring
for Seq-NMS. The classes are sorted in descendingorder by Seq-NMS
performance.

\pandocbounded{\includegraphics[keepaspectratio,alt={image}]{images/04a7bc4faf264c99a81f2358cadb8e0bcf234b77cd09e2f4a4c1613fb362d346.jpg}}

Figure 4: Absolute improvement in mAP \(( \% )\) using Seq-NMS. The
improvement is relative to singleimage NMS. Note that 7 classes have
higher than \(10 \%\) improvement, and only two classes showdecreased
performance (train and whale).

\section{5 Related Work}\label{related-work}

Many previous works in video object detection framed as multiple object
tracking. A popular sub-class of these techniques were models that did
''tracking-by-detection'', whereby a detection al-gorithm is applied on
each video frame and the detections are associated across frames to form

\pandocbounded{\includegraphics[keepaspectratio,alt={image}]{images/682b24090192d8f8c8b3dd1d29894a7402b97ca535e8195715672cd42c996713.jpg}}

Figure 5: Example video clips where Seq-NMS improves performance. The
boxes represent asequence selected by Seq-NMS. Clips are subsampled to
provide examples of high and low scoringboxes. In clips a, b, and e, the
object becomes more and more occluded as it exits the frame, leadingto
lower scores. Meanwhile, in clips c and d, the object of interest has a
low classifier score becauseit is either very small or blurred,
respectively. In all of these cases, Seq-NMS' rescoring
significantlyboosts the weaker detections by using the strong detections
from adjacent frames.

trajectories for each object. Previous detection methods were usually
based on motion {[}17{]} or ob-ject appearance {[}4{]}. With regards to
the association step, a classic method involved using Kalmanfilters to
predict tracks and the Hungarian method {[}8, 15{]} to associate
detections between frames.Particle filter techniques {[}13, 3{]} further
improved on Kalman filters by being able to handle multiplehypotheses.
Other classes of methods tried to compute all of the object trajectories
at once usinglinear programming {[}6, 2{]}. While these methods are able
to find a global optimum with high prob-ability, they assume that the
number of objects to be tracked is known a priori. On the other
hand,dynamic programming {[}12, 1{]} can also be used to find
trajectories one by one in a greedy fashion.Our proposed model is
similar in that it takes detections from a state-of-the-art single image
objectdetection method {[}9{]} and subsequently associates tracks over
time by finding the highest scoringpath by also using dynamic
programming.

\section{6 Conclusion}\label{conclusion}

By using the strong baseline of Faster R-CNN and leveraging additional
temporal information, wewere one of the top performers in the ImageNet
Object Detection from Video challenge. We wouldlike to continue pursuing
improvements to our submission, including training on the entire
VIDdataset, experimenting with neural network suppression, and
performing a deeper analysis on ourmodel designed to elucidate its
weaknesses.

\pandocbounded{\includegraphics[keepaspectratio,alt={image}]{images/ce21b8166a1f4228990f0565f67f065e75cc480f2206dfac61de6da71bea4b86.jpg}}

Figure 6: Video clips in the ImageNet VID dataset where Seq-NMS does not
improve performance.In clip a, Seq-NMS has difficulty when there are
several objects with similar appearance closetogether in the video (clip
a). This will cause the detector to drift from one object to another
whichleads to missed detections and incorrect score assignment. Seq-NMS
also accumulates spuriousdetections which leads to many false positives
(clips b and c). This occurs because Seq-NMS'objective function does not
penalize against adding more detections.

\section{Acknowledgments}\label{acknowledgments}

The six Tesla K40 GPUs used for this research were donated by the NVIDIA
Corporation.

\section{References}\label{references}

{[}1{]} Jer´ ome Bercla, Francois Fleuret, and Pascal Fua. Robust people
tracking with global trajectory opti- ˆmization. In Computer Vision and
Pattern Recognition, 2006 IEEE Computer Society Conference on,volume 1,
pages 744--750. IEEE, 2006.

{[}2{]} Jerome Berclaz, Francois Fleuret, Engin Turetken, and Pascal
Fua. Multiple object tracking us- ¨ing k-shortest paths optimization.
Pattern Analysis and Machine Intelligence, IEEE Transactions
on,33(9):1806--1819, 2011.

{[}3{]} Michael D Breitenstein, Fabian Reichlin, Bastian Leibe, Esther
Koller-Meier, and Luc Van Gool. Robusttracking-by-detection using a
detector confidence particle filter. In Computer Vision, 2009 IEEE
12thInternational Conference on, pages 1515--1522. IEEE, 2009.

{[}4{]} Pedro F Felzenszwalb, Ross B Girshick, David McAllester, and
Deva Ramanan. Object detection withdiscriminatively trained part-based
models. Pattern Analysis and Machine Intelligence, IEEE Transac-tions
on, 32(9):1627--1645, 2010.

{[}5{]} Ross Girshick. Fast R-CNN. In Proceedings of the International
Conference on Computer Vision (ICCV),2015.

{[}6{]} Hao Jiang, Sidney Fels, and James J Little. A linear programming
approach for multiple object tracking.In Computer Vision and Pattern
Recognition, 2007. CVPR'07. IEEE Conference on, pages 1--8. IEEE,2007.

{[}7{]} Kai Kang, Hongsheng Li, Junjie Yan, Xingyu Zeng, Bin Yang, Tong
Xiao, Cong Zhang, Zhe Wang,Ruohui Wang, Xiaogang Wang, et al.~T-cnn:
Tubelets with convolutional neural networks for objectdetection from
videos. arXiv preprint arXiv:1604.02532, 2016.

{[}8{]} AG Amitha Perera, Chukka Srinivas, Anthony Hoogs, Glen Brooksby,
and Wensheng Hu. Multi-objecttracking through simultaneous long
occlusions and split-merge conditions. In Computer Vision and Pat-tern
Recognition, 2006 IEEE Computer Society Conference on, volume 1, pages
666--673. IEEE, 2006.

{[}9{]} Shaoqing Ren, Kaiming He, Ross Girshick, and Jian Sun. Faster
R-CNN: Towards real-time objectdetection with region proposal networks.
In Neural Information Processing Systems (NIPS), 2015.

{[}10{]} Karen Simonyan and Andrew Zisserman. Very deep convolutional
networks for large-scale image recog-nition. arXiv preprint
arXiv:1409.1556, 2014.

{[}11{]} J.R.R. Uijlings, K.E.A. van de Sande, T. Gevers, and A.W.M.
Smeulders. Selective search for objectrecognition. International Journal
of Computer Vision, 2013.

{[}12{]} Jack K Wolf, Audrey M Viterbi, and Glenn S Dixon. Finding the
best set of k paths through a trellis withapplication to multitarget
tracking. Aerospace and Electronic Systems, IEEE Transactions on,
25(2):287--296, 1989.

{[}13{]} Changjiang Yang, Ramani Duraiswami, and Larry Davis. Fast
multiple object tracking via a hierarchicalparticle filter. In Computer
Vision, 2005. ICCV 2005. Tenth IEEE International Conference on, volume
1,pages 212--219. IEEE, 2005.

{[}14{]} Matthew D. Zeiler and Rob Fergus. Visualizing and understanding
convolutional networks. In ComputerVision - ECCV 2014 - 13th European
Conference, Zurich, Switzerland, September 6-12, 2014, Proceed-ings,
Part I, pages 818--833, 2014.

{[}15{]} Hongyi Zhang, Andreas Geiger, and Raquel Urtasun. Understanding
high-level semantics by modelingtraffic patterns. In The IEEE
International Conference on Computer Vision (ICCV), December 2013.

{[}16{]} C. Lawrence Zitnick and Piotr Dollar. Edge boxes: Locating
object proposals from edges. In ´ ECCV.European Conference on Computer
Vision, September 2014.

{[}17{]} Zoran Zivkovic. Improved adaptive gaussian mixture model for
background subtraction. In PatternRecognition, 2004. ICPR 2004.
Proceedings of the 17th International Conference on, volume 2,
pages28--31. IEEE, 2004.

\end{document}
